\section{Síntesis y ampliación}

\subsection{Conclusiones centrales}

\begin{enumerate}
\item La trayectoria personal de Xie Saining (谢赛宁) no es un manual de éxito, sino una muestra de cómo un problema de investigación toma forma de manera gradual mediante elecciones, organización y azar.
\item La visión no debe entenderse como un conjunto de tareas antiguas, sino como una perspective para procesar señales del mundo real continuas, de alta dimensión y con ruido.
\item La definición mínima de un modelo del mundo es $s_{t+1}=f(s_t,a_t)$, pero lo verdaderamente difícil es la representación del state, el anclaje físico, el planning y el ciclo cerrado de retroalimentación.
\item El LLM es un componente importante de un sistema inteligente, pero se alimenta sobre todo del conocimiento que la humanidad ya ha tokenizado y no puede sustituir de forma automática el modelado del mundo físico.
\item «OpenAI a la inversa» significa que un modelo del mundo no puede limitarse a descargar internet, sino que debe construir junto con socios del mundo real ciclos cerrados de datos, evaluación y aplicación.
\item Lo esencial de AMI Labs no es ser «otra empresa de modelos grandes», sino intentar una nueva forma de organización que equilibre libertad de investigación, escala de recursos, ciclo comercial cerrado y vínculo con la academia.
\item El problema de estar LLM-pilled es que una narrativa única define las tablas de clasificación y la asignación de recursos, con lo que reduce el espacio de exploración de direcciones como los modelos del mundo, la comprensión de video y la inteligencia física.
\item La inteligencia no es una sola línea de puntuación lingüística; la inteligencia del mundo real tiene que volver al cuerpo, el entorno, la acción, la retroalimentación y los objetivos.
\end{enumerate}

\subsection{Preguntas abiertas}

Después de la síntesis conviene mantener una actitud de interrogación, porque uno de los aspectos más honestos de este episodio es reconocer que muchas respuestas clave sobre los modelos del mundo todavía no han convergido. Las siguientes preguntas son también pistas que deben seguirse de manera continua en lecturas posteriores y al generar nuevas notas.

\begin{itemize}
\item ¿El preentrenamiento de modelos del mundo dará lugar a un paradigma estable similar a la next-token prediction?
\item ¿El state de un modelo del mundo debe ser 3D explícito, una latent representation, una estructura híbrida o formarse de manera dinámica según la tarea?
\item ¿La frontera entre VLA y world model se irá fusionando o se formará una división del trabajo entre una capa base aguas arriba y una capa de acción aguas abajo?
\item ¿Una red global de socios al estilo de AMI podrá aportar realmente un ciclo cerrado de datos insustituible?
\item Cuando las tablas de clasificación de LLM se saturen, ¿la industria volverá a dirigir recursos hacia la physical intelligence?
\end{itemize}

\subsection{Lecturas adicionales}

\begin{itemize}
\item Yann LeCun, \textit{A Path Towards Autonomous Machine Intelligence}: para entender la ruta de largo plazo de JEPA, los modelos del mundo y la autonomous intelligence.
\item Richard Sutton, artículos sobre Dyna / model-based reinforcement learning: para entender la relación histórica entre los modelos del mundo, el planning y el RL.
\item Materiales introductorios de Model Predictive Control: para entender la intuición de control de «usar un modelo para hacer roll out del futuro y elegir una acción».
\item Artículos sobre aprendizaje de representaciones visuales, generación de video, 3D representation, VLA y preentrenamiento de robots: para entender las múltiples vías de entrada posibles hacia los modelos del mundo.
\item Frans de Waal, \textit{Are We Smart Enough to Know How Smart Animals Are?}: para entender una visión de la inteligencia no antropocéntrica.
\end{itemize}

\begin{importantbox}{Juicio final}
Lo que más vale la pena llevarse de este episodio no es la dicotomía «los LLM están equivocados y los modelos del mundo tienen razón», sino un juicio más preciso: los LLM han resuelto una gran parte del mundo que la humanidad ya puso por escrito; los modelos del mundo deben resolver el mundo que la humanidad no ha escrito y que es muy difícil de escribir. Lo decisivo para la siguiente etapa de la inteligencia quizá se encuentre en el puente entre ambos.
\end{importantbox}

\end{document}
